% ============================================================
% PRN232 - Section 9: Conclusion (Vietnamese)
% ============================================================
\section{\bi{Kết luận và hướng phát triển}{Conclusion and Future Work}}
\label{sec:concl}

\subsection{\bi{Kết luận}{Conclusion}}

\bi{
Bài báo đã trình bày OmniCast --- một hệ thống lịch phát và phát sóng
nội dung số đa nền tảng xây dựng cho môn PRN232. Ba đóng góp công
nghệ chính được chứng minh bằng mã nguồn và kiểm thử thực tế: thuật
toán phát hiện xung đột lịch phát bằng một truy vấn SQL thô tính
thời điểm kết thúc hiệu lực; cơ chế điền lịch EPG 24/7 với lựa chọn
ngẫu nhiên có hạt giống bảo đảm tính tái lập; và quy trình biên tập
AI gọi mô hình một lần duy nhất, kiểm tra schema bằng Zod, có đường
dự phòng heuristic và dấu vết phiên bản đầy đủ.

Toàn bộ hệ thống --- 19 module NestJS, 140 endpoint REST, 21 model
Prisma, cổng web Next.js 15 với 33 route và microservice ghi nhật ký
kiểm toán --- đã được triển khai trên Vercel, kiểm thử bằng 482 kiểm
thử tự động trong quy trình tích hợp liên tục, và vận hành với ba
điểm cuối sản xuất trả về HTTP 200.
}{
This paper has presented OmniCast, a multi-platform digital broadcast
scheduling and delivery system built for PRN232. Its three principal
technical contributions are demonstrated through the source code and
real test runs: a broadcast-conflict algorithm implemented as a single
raw SQL query that computes the effective end time; a 24/7 EPG filler
whose seeded selection guarantees reproducibility; and an AI
curation pipeline that issues a single consolidated model call,
validates the output against a Zod schema, offers a heuristic fallback
and records full version provenance.

The complete system --- 19 NestJS modules, 140 REST endpoints, 21
Prisma models, a Next.js 15 portal with 33 routes and a gRPC audit
microservice --- has been deployed on Vercel, validated by 482
automated tests inside a continuous-integration pipeline, and operates
with three production endpoints returning HTTP 200.
}

\subsection{\bi{Hạn chế đã biết}{Known Limitations}}

\bi{
Ba hạn chế đã biết cần được nêu rõ. Thứ nhất, mức độ phủ kiểm thử
E2E còn thấp: tám kịch bản trình duyệt không bao gồm luồng phát sóng
trực tiếp vốn cần phần cứng và cơ sở hạ tầng MediaMTX. Thứ hai, lớp
ingest hiện hỗ trợ YouTube, Twitch và RSS nhưng chưa hỗ trợ nguồn
tương tác thời gian thực. Thứ ba, thuật toán điền lịch chọn theo thể
loại và theo hạt giống, chưa tối ưu theo hành vi người xem thực tế.
}{
Three known limitations should be stated. First, end-to-end coverage
remains low: eight browser scenarios exclude the live-broadcast flow,
which requires dedicated hardware and MediaMTX infrastructure.
Second, the ingestion layer supports YouTube, Twitch and RSS but not
real-time interactive sources. Third, the filler algorithm selects by
category and by seed rather than by observed viewer behaviour.
}

\subsection{\bi{Hướng phát triển}{Future Work}}

\bi{
Bốn hướng phát triển được ưu tiên cho các phiên bản tiếp theo.
}{
Four directions are prioritised for subsequent iterations.
}

\begin{enumerate}[leftmargin=*,itemsep=1pt]
  \item \bi{\textbf{Cá nhân hoá lịch phát bằng học máy.} Thay vì
  chọn theo thể loại, thuật toán điền lịch có thể dùng lịch sử xem
  và dữ liệu tương tác để đề xuất nội dung cho từng kênh, đồng thời
  giữ nguyên tính chất xác định để có thể kiểm thử.}{
  \textbf{Recommender-driven scheduling.} Instead of selecting by
  category, the filler algorithm could use watch history and
  interaction data to propose content per channel while preserving
  determinism so the result remains testable.}

  \item \bi{\textbf{Mở rộng kiểm thử E2E.} Xây dựng môi trường
  MediaMTX trong CI để kiểm thử tự động luồng WHIP phát và HLS phát
  lại, đồng thời bổ sung kiểm thử trình duyệt cho vai trò STAFF và
  ADMIN.}{
  \textbf{Scaling end-to-end testing.} Provision a MediaMTX instance
  in CI to test the WHIP publish and HLS playback paths
  automatically, and add browser scenarios for the STAFF and ADMIN
  roles.}

  \item \bi{\textbf{Cải thiện trợ lý AI.} Bổ sung đánh giá nhiều mô
  hình, kiểm duyệt đầu ra bằng mô hình thứ hai và xếp hạng các báo
  cáo AI theo phản hồi của biên tập viên để tinh chỉnh prompt theo
  vòng lặp.}{
  \textbf{Improving the AI assistant.} Add multi-model evaluation,
  second-model output verification, and ranking of AI reports by
  editor feedback so the prompt can be tuned iteratively.}

  \item \bi{\textbf{Tách lớp truy cập dữ liệu.} Đưa truy vấn EPG vào
  một kho phân tích chỉ đọc với các chỉ mục bao phủ để tách tải
  đọc khỏi tải ghi trên cùng một cơ sở dữ liệu.}{
  \textbf{Separating the data-access tier.} Move EPG queries into a
  read-only analytical store with covering indexes to separate read
  load from write load on the same database.}
\end{enumerate}

\subsection{\bi{Tính xác thực và khả năng tái lập}{Reproducibility Statement}}

\bi{
Toàn bộ kết quả trong bài báo có thể tái lập. Mã nguồn nằm trong kho
công khai \texttt{TomOutfit/OmniCast\_PRN232-PRM393\_Project\_Fa26Semester};
ba lệnh chạy kiểm thử là \texttt{npm test} trong thư mục
\texttt{backend}, \texttt{npx vitest run} trong \texttt{frontend} và
\texttt{flutter test} trong \texttt{mobile}. Bảng kết quả nguyên bản
được đính kèm trong tài liệu kiểm thử đi kèm.
}{
Every result in this paper is reproducible. The source code lives in
the public repository
\texttt{TomOutfit/OmniCast\_PRN232-PRM393\_Project\_Fa26Semester}; the
three test commands are \texttt{npm test} in \texttt{backend},
\texttt{npx vitest run} in \texttt{frontend} and \texttt{flutter test}
in \texttt{mobile}. The raw result files are attached to the
accompanying testing documentation.
}
